\section{Proofs for abelian categories and exact sequences}
\label{annex:p10-abelian-exact}

\begin{annexitem}{annex:p10.ab.1}{thm:p10-ab-normal}{Normal universal arrows and finite constructions}
\begin{proof}
Kernel arrows are monic: if \(kr=ks\), both \(r,s\) are the unique
factorization of that arrow through the kernel. Cokernel arrows are epic
by the dual uniqueness argument. If \(1_T=0\), every arrow into or out
of \(T\) is zero, so \(T\) is a zero object.

Let a mono \(m:M\to A\) be a kernel of \(h:A\to D\), and put
\(q=\operatorname{Cok}m\), \(k=\operatorname{Ker}q\).
Since \(hm=0\), \(h=tq\) for a unique \(t\). Since \(qm=0\),
\(m=ku\); and since \(hk=tqk=0\), \(k=mv\).
Cancellation of the two monos in these equations gives \(uv=1\) and
\(vu=1\). Thus \(m\) is, compatibly, the kernel of its own cokernel.
Applying this argument in the opposite category proves that each epi is
the cokernel of its own kernel. If \(m\) is also epic, \(hm=0\)
forces \(h=0\); the kernel of this zero arrow is \(1_A\), so \(m\)
is an isomorphism by uniqueness of kernels.

For monos \(m:M\to A,n:N\to A\), take \(q=\operatorname{Cok}m\)
and \(s:S\to N=\operatorname{Ker}(qn)\). Since \(qns=0\), there
is a unique \(t:S\to M\) with \(mt=ns\). For any \(x:T\to M\),
\(y:T\to N\) with \(mx=ny\), \(qny=0\), so \(y=sz\) for
unique \(z\). Then \(mx=nsz=mtz\), hence \(x=tz\). This proves
the pullback property and the greatest-lower-subobject property; \(ns\)
is monic. The dual construction supplies intersections of quotient
objects. The kernel/cokernel inverse correspondence just proved reverses
containment, so the latter construction supplies binary unions of
subobjects as well.

For parallel \(f,g:A\to B\), the graphs
\(\langle1,f\rangle,\langle1,g\rangle:A\to A\times B\) are
monic. Take their pullback intersection, with maps \(u,v:E\to A\).
Its graph equation gives \(u=v=e\) after the first projection and
\(fe=ge\) after the second. If \(fx=gx\), the two graph maps of
\(x:T\to A\) coincide, so the pullback gives a unique \(z:T\to E\)
with \(ez=x\). Thus \(e\) is an equalizer. Duality supplies all
coequalizers. With zero objects and binary products/coproducts,
Theorem~\ref{thm:p10-limits-products-equalizers} now constructs all
finite limits/colimits, including arbitrary pullbacks/pushouts. No hom
addition has been used.
\end{proof}
\end{annexitem}

\begin{annexitem}{annex:p10.ab.2}{thm:p10-ab-factorization}{Images and the non-assumed coimage comparison}
\begin{proof}
Let \(q:B\to Q=\operatorname{Cok}f\) and
\(i:I\to B=\operatorname{Ker}q\). The equation \(qf=0\) gives
\(f=ie\) for unique \(e:A\to I\). If \(f=mx\) for any mono
\(m:M\to B\), its cokernel \(q_m\) kills \(f\), hence
\(q_m=tq\). Thus \(q_m i=0\), and \(i\) factors through
\(m=\operatorname{Ker}q_m\). This proves the least-subobject property.
The categorical dual proves the stated coimage formula.

If \(\operatorname{Cok}f=0\), its image is \(1_B\). For arrows
\(r,s:B\to D\) with \(rf=sf\), their equalizer allows \(f\),
and therefore allows its image \(1_B\). The equalizer is consequently
invertible, giving \(r=s\). Thus \(f\) is epic. Conversely, if
\(f\) is epic, every arrow killing it is zero; \(B\to0\) has the
cokernel universal property. The dual proves
\(f\) monic iff \(\operatorname{Ker}f=0\).

To prove that \(e:A\to I\) above is epic, let its image be
\(j:J\to I\), and write \(e=jd\). Then \(ij\) is a subobject
allowing \(f\). Leastness of \(i\) gives \(i=ijh\), so
\(jh=1_I\). As \(j\) is monic, \(hj=1_J\); hence \(j\) is
invertible and \(\operatorname{Im}e=I\). The proved epicity criterion
now makes \(e\) epic. Because \(i\) is monic, the universal kernel
of \(e\) is also the kernel of \(f\): an arrow is killed by one
iff it is killed by the other. Normality therefore says that \(e\)
is a cokernel of \(\operatorname{Ker}f\).

Both \(q_f:A\to\operatorname{Coim}f\) and \(e:A\to I\) are
cokernels of this same kernel. Their universal properties give inverse
compatible maps \(\bar f:\operatorname{Coim}f\to I\) and
\(\bar f^{-1}:I\to\operatorname{Coim}f\). This establishes, rather
than assumes, invertibility of the comparison.
For any factorization \(f=me'\) with \(e'\) epic and \(m\) monic,
\(\operatorname{Ker}e'=\operatorname{Ker}f\) and
\(\operatorname{Cok}m=\operatorname{Cok}f\), by cancellation and the
two universal properties. Normality identifies its two arrows with the
same coimage and image. The unique comparison maps are inverse because
their composites have the identity after epic precomposition by \(e'\)
or \(q_f\). This proves the promised uniqueness as well.
\end{proof}
\end{annexitem}

\begin{annexitem}{annex:p10.ab.3}{thm:p10-ab-biproduct}{The canonical sum--product map}
\begin{proof}
For the coproduct \(S=A+B\), let \(u_A,u_B\) be its injections and
\(r_A=[1,0],r_B=[0,1]\). They satisfy the four projection equations.
The injection \(u_A\) is monic because \(r_Au_A=1\).
The map \(r_B\) is a cokernel of \(u_A\): a map \([x,y]:S\to T\)
kills \(u_A\) iff \(x=0\), in which case it equals \(yr_B\),
uniquely. Normality gives \(u_A=\operatorname{Ker}r_B\).
The corresponding assertions hold with \(A,B\) interchanged.

For the product \(P=A\times B\), its injections
\(j_A=\langle1,0\rangle,j_B=\langle0,1\rangle\) satisfy
\(j_A=\operatorname{Ker}p_B\). The split epi \(p_B\) is therefore
its cokernel by normality. Let \(\sigma:S\to P\) be the canonical
map, so \(p_A\sigma=r_A,p_B\sigma=r_B\) and
\(\sigma u_A=j_A,\sigma u_B=j_B\).
If \(\sigma x=0\), then \(r_Bx=0\), so \(x=u_Ay\);
\(r_Ax=y=0\). Hence \(\sigma\) has zero kernel and is monic.
If \(z\sigma=0\), then \(zj_A=0\), hence \(z=wp_B\);
\(zj_B=w=0\). Thus its cokernel is zero and it is epic.
Balance makes \(\sigma\) invertible. This proof never adds arrows.
The transported product structure makes \(S=A\oplus B\) a biproduct.
Finite iteration gives all finite biproducts; the empty one is \(0\).
\end{proof}
\end{annexitem}

\begin{annexitem}{annex:p10.ab.4}{thm:p10-ab-addition}{Interchange and the shear that constructs a negative}
\begin{proof}
Define two initially distinct operations on arrows \(A\to B\):
\(L(f,g)=[f,g]\Delta_A\) and
\(R(f,g)=\nabla_B\langle f,g\rangle\).
The biproduct equations give a common zero unit on both sides: for
example \([f,0]=fp_1\) and \(p_1\Delta_A=1\), so \(L(f,0)=f\);
the other three unit equations follow from the other projection/injection.
Postcomposition distributes over \(L\) by the coproduct universal
property, and precomposition over \(R\) by the product universal property.

Given four arrows \(a,b,c,d:A\to B\), take the unique map
\(M:A\oplus A\to B\oplus B\) with these four components.
Computing \(\nabla_BM\Delta_A\) first by product components and
then by coproduct components gives the interchange equation
\[
 R(L(a,b),L(c,d))=L(R(a,c),R(b,d)).
\]
Set \(b=c=0\) to obtain \(R(a,d)=L(a,d)\); call the common
operation \(+\). Interchange becomes
\((a+b)+(c+d)=(a+c)+(b+d)\). Taking \(c=0\) proves associativity,
and taking \(a=d=0\) proves commutativity. Both distributive
composition laws now hold. This is a commutative monoid structure before
any inverse has been asserted. Its expression is precisely
\(\nabla_B(f\oplus g)\Delta_A\).

For clarity we use Freyd's row-source matrix convention here: the
\((i,j)\)-entry of \(M:\bigoplus_i A_i\to\bigoplus_j B_j\) is
\(p_jMu_i:A_i\to B_j\). Let
\(S=\left(\begin{smallmatrix}1&f\\0&1\end{smallmatrix}\right)
:A\oplus B\to A\oplus B\), where \(f:A\to B\).
Its product components on an arrow \(x:T\to A\oplus B\) are
\(p_Ax\) and \(fp_Ax+p_Bx\). If \(Sx=0\), both components of
\(x\) are zero, so \(x=0\). Thus \(S\) is monic.
For \(z:A\oplus B\to T\), \(zS=0\) gives
\(zu_B=0\) and \(zu_A+zu_Bf=0\), hence \(zu_A=0\) and \(z=0\).
So \(S\) is epic, and therefore invertible. Write its inverse with
four components. The equation \(SS^{-1}=1\), composed with the
projections and injections, gives the diagonal entries \(1,1\), the
lower-left entry zero, and an upper-right entry \(g:A\to B\) with
\(f+g=0\). This constructs an inverse of every \(f\). The
commutative monoids are thus abelian groups.

Finally, in any given bilinear hom-group structure, the four biproduct
equations imply \(u_1p_1+u_2p_2=1\), since their composites with
both product projections equal those of the identity. They then express
the given sum by the same diagonal/codiagonal formula. Hence the
canonical addition is unique. The entire argument is categorical;
no underlying elements or embedding theorem enter it.
\end{proof}
\end{annexitem}

\begin{annexitem}{annex:p10.ab.5}{thm:p10-ab-additive-bridge}{Biproduct equations and the two presentations}
\begin{proof}
In a preadditive category with zero object and a product \(P=A\times B\),
form \(u_1=\langle1,0\rangle,u_2=\langle0,1\rangle\).
The composites of \(u_1p_1+u_2p_2\) with both projections are those
of \(1_P\), so they are equal. For arrows \(f:A\to X,g:B\to X\),
put \(h=fp_1+gp_2\). Then \(hu_1=f,hu_2=g\); conversely any
such \(h\) equals \(h(u_1p_1+u_2p_2)=fp_1+gp_2\).
Thus the product is a coproduct. The same calculation proves that the
two displayed systems of equations suffice for a biproduct; necessity
comes from its two universal properties and bilinearity. Finite
iteration, including the zero object, proves the finite assertion.

For maps represented by entries \(f_{ij}:A_i\to B_j\) and
\(g_{jk}:B_j\to C_k\), insert
\(1_{\bigoplus B_j}=\sum_j u_jp_j\) into the composite. Its
\((i,k)\)-entry is \(\sum_j g_{jk}f_{ij}\). Empty sums give zero.
The structural axioms yield a preadditive category by the preceding
proof; conversely the additive presentation supplies products and,
by this construction, coproducts, and its remaining hypotheses are
exactly the kernels/cokernels/normality axioms. Uniqueness of the
canonical addition identifies the two structures, completing both
directions of the bridge.
\end{proof}
\end{annexitem}

\begin{annexitem}{annex:p10.ab.6}{thm:p10-ab-exactness}{Exactness without unspecified elements}
\begin{proof}
Write \(i_f:\operatorname{Im}f\to B=\operatorname{Ker}q_f\).
If \(gf=0\), \(g\) factors through \(q_f\), so \(g i_f=0\)
and \(\operatorname{Im}f\subseteq\operatorname{Ker}g\).
If also \(q_f k_g=0\), \(k_g\) factors through \(i_f\), giving
the opposite containment. Conversely equality of these subobjects
gives both equations by their universal properties.
Applying the inverse kernel/cokernel correspondence to the equality
gives \(\operatorname{Cok}(\operatorname{Ker}g)=\operatorname{Cok}f\),
since \(f=i_fe_f\) with \(e_f\) epic has the same cokernel as \(i_f\).
This is the equivalent quotient-object assertion.

At the left endpoint, the image of the zero arrow is zero; hence
exactness of \(0\to A\xrightarrow f B\) says \(\operatorname{Ker}f=0\),
equivalent to monicity. At the right endpoint, the kernel of \(B\to0\)
is \(1_B\), so exactness says \(\operatorname{Im}f=B\), equivalent
to epicity. Combining both endpoints gives balance and the isomorphism
case. In a short exact sequence, the mono \(i\) identifies its source
with its image, hence is a kernel of \(p\); the epi \(p\) is then a
cokernel of \(i\) by normality. Conversely these two universal
properties give both endpoint conditions and exactness in the middle.
For any \(f=i_fe_f\), \(e_f\) is epic with kernel \(\operatorname{Ker}f\),
and \(i_f\) is monic with cokernel \(\operatorname{Cok}f\); this
constructs its two asserted short sequences.
\end{proof}
\end{annexitem}

\begin{annexitem}{annex:p10.ab.7}{thm:p10-ab-splitting}{Actual inverse maps for a split extension}
\begin{proof}
Suppose \(ps=1_C\). Since \(p(1_B-sp)=0\), the kernel \(i\)
gives unique \(r:B\to A\) with \(ir=1_B-sp\).
Composing with \(i\) gives \(iri=i\), so \(ri=1_A\);
composing with \(s\) gives \(irs=0\), so \(rs=0\).
Together with \(pi=0,ps=1\), the equation \(ir+sp=1_B\)
shows that \((i,s):A\oplus C\to B\) and
\(\langle r,p\rangle:B\to A\oplus C\) are inverse maps.

Conversely, if \(ri=1_A\), then \((1_B-ir)i=0\), so the cokernel
\(p\) gives unique \(s:C\to B\) with \(sp=1_B-ir\).
Composing with \(p\) and cancelling the epic \(p\) yields \(ps=1_C\).
The same inverse-map calculation follows. A compatible biproduct
isomorphism supplies these section and retraction maps via its injection
and projection. This proves all three implications and not just an
object-level isomorphism.
\end{proof}
\end{annexitem}

\begin{annexitem}{annex:p10.ab.8}{thm:p10-ab-base-change}{Signed sequences and epic base change}
\begin{proof}
Let \(h=(f,-g):A\oplus B\to D\), with kernel \(k:P\to A\oplus B\).
For a commuting pair \(x:T\to A,y:T\to B\), the pairing
\(\langle x,y\rangle\) is killed by \(h\); its unique kernel
factorization gives a unique map to \(P\) with both prescribed
projections. This is the pullback. The signed sequence is therefore
left exact precisely when the given square has that universal property.
In the dual category it is right exact precisely when the square is a
pushout. Equivalently, a compatible cocone \(x:A\to T,y:B\to T\)
corresponds to the row \((x,-y)\) killing \((a,b):P\to A\oplus B\),
so the signed cokernel universal property is exactly the pushout
property. Both conditions together give the displayed short sequence.

Suppose \(g\) is epic. The map \(h\) is epic, because its restriction
to \(B\) is \(-g\), an epi. Put \(r=p_Ak:P\to A\), and
let \(q:A\to Q\) be its cokernel. The row \((q,0):A\oplus B\to Q\)
kills \(k\), so factors through \(h=\operatorname{Cok}k\) as
\((q,0)=th\). Restricting to \(B\) gives \(0=-tg\), so \(t=0\)
and \(q=0\). Thus \(r\) is epic.
If \(z:T\to P\) is killed by \(r\), its other projection is killed
by \(g\), and gives the unique factorization through
\(\operatorname{Ker}g\). Conversely, the pair consisting of a zero
arrow to \(A\) and the kernel arrow to \(B\) has a unique map into
\(P\). These constructions are inverse universal kernel properties,
so the kernel of \(r\) is the pulled-back kernel of \(g\).
They also show, for arbitrary \(g\), the kernel identification itself.
The opposite-category argument proves the claimed pushout stability of
monos and their cokernels, with no assumption about modules or elements.
\end{proof}
\end{annexitem}

\begin{annexitem}{annex:p10.ab.9}{thm:p10-ab-local-arrows}{Epic local lifts justify the later chase}
\begin{proof}
Factor \(f=ie\) with \(e:A\to I\) epic and \(i:I\to B\) monic.
If \(x=i\bar x\), pull \(e\) back along \(\bar x:T\to I\).
The resulting \(s:E\to T\) is epic by base change, with an arrow
\(y:E\to A\) satisfying \(ey=\bar xs\), hence \(fy=xs\).
Conversely, if \(fy=xs\) with \(s\) epic, then
\(q_fx s=q_ffy=0\), so \(q_fx=0\); the kernel \(i\) supplies
the factorization of \(x\). Equality after an epic precomposition is
exactly the defining cancellation property of an epi. The composite
of two epis is epic, by two successive cancellations.

For exactness testing, suppose \(hr=0\). This gives
\(\operatorname{Im}r\subseteq\operatorname{Ker}h\).
If the universal kernel arrow \(k:K\to B\) of \(h\) admits an
epic local lift through \(r\), the just-proved criterion factors \(k\)
through \(\operatorname{Im}r\), proving the opposite containment.
Thus showing that each arrow killed by \(h\) locally lifts through
\(r\) is sufficient, since it includes this particular \(k\).
Every local refinement below is a finite composition of explicitly
constructed pullbacks of epis. There is no assumed set-valued
interpretation, no choice of elements at every object, and no unproved
embedding theorem behind this argument.
\end{proof}
\end{annexitem}

\begin{annexitem}{annex:p10.ab.10}{thm:p10-ab-short-five}{Separate kernel and cokernel arguments}
\begin{proof}
Use row maps \(i,p\) and \(i',p'\), so \(bi=i'a\) and
\(p'b=cp\). If \(x:T\to B\) satisfies \(bx=0\), then
\(cpx=p'bx=0\). Since \(c\) is monic, \(px=0\), giving
\(x=iy\). Now \(i'ay=biy=0\), so \(y=0\) because both
\(i'\) and \(a\) are monic. Therefore \(x=0\), and the zero-kernel
criterion makes \(b\) monic.

Let \(q:B'\to Q\) kill \(b\). Then \(qi'a=qbi=0\);
the epic \(a\) implies \(qi'=0\). Since \(p'\) is a cokernel
of \(i'\), \(q=tp'\). The equation \(qb=0\) now reads
\(tcp=0\). Cancelling the epis \(p,c\) gives \(t=0\) and
\(q=0\). Thus \(b\) has zero cokernel and is epic. Balance gives
an actual inverse of \(b\), completing the lemma.
\end{proof}
\end{annexitem}

\begin{annexitem}{annex:p10.ab.11}{thm:p10-ab-snake}{Constructing the connecting arrow and proving all positions}
\begin{proof}
Write the two rows as \((i,p)\), \((i',p')\), their vertical arrows
as \(a,b,c\), and the kernel arrows as \(k_a,k_b,k_c\), with
cokernel arrows \(q_a,q_b,q_c\). The commutativity equations are
\(bi=i'a\) and \(p'b=cp\).

Form the pullback \(P=B\times_C\operatorname{Ker}c\) with
\(u:P\to B,v:P\to\operatorname{Ker}c\). Since \(p\) is epic,
\(v\) is epic. The pair \((i,0)\) gives \(j:A\to P\), and
\(j\) is a kernel of \(v\): if \(vx=0\), then \(pux=0\),
so \(ux=iy\) uniquely, and the pullback property gives \(x=jy\).
Thus \(v\) is a cokernel of \(j\). Since
\(p'bu=cpu=ck_cv=0\), there is unique \(t:P\to A'\) with
\(i't=bu\). The equality \(i'tj=bi=i'a\) gives \(tj=a\).
Therefore \(q_atj=0\), and the cokernel \(v\) gives unique
\(\delta:\operatorname{Ker}c\to\operatorname{Cok}a\) with
\(\delta v=q_at\). This establishes its construction and sign.

The induced arrows \(\alpha,\beta,\gamma,\eta\) between the other
terms are uniquely defined by
\[
 k_b\alpha=ik_a,\quad k_c\beta=pk_b,\quad
 \gamma q_a=q_bi',\quad \eta q_b=q_cp'.
\]
Their existence follows respectively from the kernel and cokernel
equations of the commutative diagram. The composite \(ik_a\) is monic,
so \(\alpha\) is monic; \(q_cp'=\eta q_b\) is epic, so
\(\eta\) is epic. These prove exactness at the endpoint zeros.

At \(\operatorname{Ker}b\), \(\beta\alpha=0\). If
\(x:T\to\operatorname{Ker}b\) satisfies \(\beta x=0\), then
\(pk_bx=0\), so \(k_bx=iy\) uniquely. Now
\(i'ay=biy=bk_bx=0\), hence \(ay=0\) and \(y=k_az\).
Cancellation of \(k_b\) gives \(x=\alpha z\). This proves exactness
there by the kernel universal property.

At \(\operatorname{Ker}c\), the pair \((k_b,\beta)\) gives a
map \(\ell:\operatorname{Ker}b\to P\). Since \(i't\ell=bk_b=0\),
\(t\ell=0\); consequently \(\delta\beta=0\).
If \(x:T\to\operatorname{Ker}c\) is killed by \(\delta\), pull
\(v\) back along \(x\). This gives an epic cover \(s:E\to T\)
and \(z:E\to P\) with \(vz=xs\). Then \(q_atz=0\).
The generalized-arrow lemma gives a further epic cover \(r:F\to E\)
and \(y:F\to A\) with \(ay=tzr\). Define
\(w=uzr-iy:F\to B\). Its image under \(b\) is zero, so
\(w=k_bw'\) for unique \(w':F\to\operatorname{Ker}b\).
Moreover \(pw=k_cxsr\), giving \(\beta w'=xsr\) after cancelling
\(k_c\). Thus every such \(x\) locally lifts through \(\beta\);
the generalized-arrow criterion proves exactness at this position.

At \(\operatorname{Cok}a\),
\(\gamma\delta v=q_bi't=q_bbu=0\), hence \(\gamma\delta=0\).
For \(x:T\to\operatorname{Cok}a\) with \(\gamma x=0\), pull
back the epic \(q_a\) to obtain \(s:E\to T,z:E\to A'\) with
\(q_az=xs\). Then \(q_bi'z=0\), so after a further epic
\(r:F\to E\) there is \(y:F\to B\) with \(by=i'zr\).
The equality \(cpy=p'by=0\) gives \(py=k_cw\) for some
\(w:F\to\operatorname{Ker}c\). The pair \((y,w)\) gives
\(h:F\to P\). Since \(i'th=by=i'zr\), \(th=zr\).
Thus \(\delta w=q_ath=q_azr=xsr\), a local lift through
\(\delta\). The criterion proves exactness here.

At \(\operatorname{Cok}b\), \(\eta\gamma=0\) follows after
precomposition by the epic \(q_a\) from \(p'i'=0\).
For \(x:T\to\operatorname{Cok}b\) with \(\eta x=0\), lift
through \(q_b\) after \(s:E\to T\), obtaining \(z:E\to B'\)
with \(q_bz=xs\). Since \(q_cp'z=0\), after an epic
\(r:F\to E\) obtain \(y:F\to C\) with \(cy=p'zr\).
Pull back the epic \(p\) along \(y\); after an epic \(d:G\to F\)
there is \(w:G\to B\) with \(pw=yd\). The arrow
\(zrd-bw:G\to B'\) is killed by \(p'\), so equals \(i'h\)
for some \(h:G\to A'\). Consequently
\(\gamma q_ah=q_bi'h=q_bzrd=xsrd\).
This is a local lift through \(\gamma\), proving the remaining
interior exactness. All covers used are epis supplied by pullbacks, and
their finite composites are epis; these are categorical arrows, not
elements of unspecified underlying modules.

For naturality, take a morphism between two such diagrams, with row
maps \(x_A,x_B,x_C\) and \(x_{A'},x_{B'},x_{C'}\).
The universal kernels/cokernels induce maps \(h_K\) between the
\(\operatorname{Ker}c\)'s and \(h_Q\) between the
\(\operatorname{Cok}a\)'s. The pullback gives \(m:P\to\widetilde P\)
with \(\widetilde u m=x_Bu\), \(\widetilde v m=h_Kv\).
The commuting squares and monicity of \(\widetilde i'\) imply
\(\widetilde t m=x_{A'}t\). Therefore
\[
 \widetilde\delta h_Kv
 =\widetilde q_a\widetilde t m
 =\widetilde q_a x_{A'}t=h_Qq_at=h_Q\delta v.
\]
Cancel the epic \(v\) to obtain
\(\widetilde\delta h_K=h_Q\delta\). This proves naturality for
the actual constructed connecting map, not just existence of some map
making a sequence exact.
\end{proof}
\end{annexitem}

\begin{annexitem}{annex:p10.ab.12}{thm:p10-ab-nine}{Both directions of the three-by-three lemma}
\begin{proof}
Call the column objects \(A_j,B_j,C_j\), \(j=1,2,3\), with
column short sequences \(0\to A_j\to B_j\to C_j\to0\).
Suppose first the top and middle rows are short exact. Apply the snake
lemma to the diagram from the top row to the middle row. Its three
vertical maps are monic, so their kernels are zero; their cokernels
are compatibly \(C_1,C_2,C_3\). The snake sequence reduces to
\(0\to C_1\to C_2\to C_3\to0\), making the bottom row short
exact. The induced maps are the given bottom maps by cokernel
uniqueness and commutativity of the columns.

Conversely, if the middle and bottom rows are short exact, apply the
snake lemma to their diagram. The vertical maps are epic, so their
cokernels are zero, while their kernels are compatibly \(A_1,A_2,A_3\).
The sequence reduces to \(0\to A_1\to A_2\to A_3\to0\), the
given top row by kernel uniqueness. This proves both directions under
the stated column and middle-row hypotheses.
\end{proof}
\end{annexitem}

\begin{annexitem}{annex:p10.ab.13}{thm:p10-ab-noether}{The categorical subquotient constructions}
\begin{proof}
For \(A\subseteq B\subseteq C\), let
\(q_A:C\to C/A\), \(q_B:C\to C/B\) be cokernels of their
inclusions. The inclusion \(B\to C\) followed by \(q_A\) has kernel
the given \(A\to B\), by the pullback kernel property. Its canonical
factorization therefore gives a mono \(u:B/A\to C/A\).
Because \(q_B\) kills \(A\), it factors as \(v q_A\); the map
\(v:C/A\to C/B\) is epic, since \(q_B\) is epic.
We have \(vu=0\). If \(x:T\to C/A\) is killed by \(v\), lift
it through \(q_A\) after an epic cover \(s:E\to T\), as
\(q_Ay=xs\). Now \(q_By=0\), so \(y\) factors through \(B\).
Thus \(xs\) factors through \(u\). The local-arrow criterion gives
\(\operatorname{Ker}v=\operatorname{Im}u\). This constructs the
short exact sequence; its last map is the cokernel of \(u\), yielding
\((C/A)/(B/A)\cong C/B\) with compatible arrows.

For \(U,V\subseteq X\), factor
\([i_U,i_V]:U\oplus V\to X\) as an epi to \(S=U+V\), followed
by a mono. Its restrictions give monos \(U\to S,V\to S\).
Let \(q:S\to S/V\) be the latter's cokernel, and let
\(h:U\to S/V\) be its composite with \(U\to S\).
The kernel of \(h\) is \(U\times_SV\), by pullback of the kernel
\(V\to S\) of \(q\); this is also \(U\times_XV\), since
\(S\to X\) is monic.
If \(z:S/V\to T\) kills \(h\), then \(zq\) kills both \(U\)
and \(V\), hence kills the epi \(U\oplus V\to S\). Thus
\(zq=0\) and \(z=0\); the cokernel criterion makes \(h\) epic.
Its coimage--image comparison consequently gives
\(U/(U\cap V)\cong S/V\). This supplies actual quotient objects
and maps, not a set-element identification.
\end{proof}
\end{annexitem}

\begin{annexitem}{annex:p10.ab.14}{thm:p10-ab-functors}{What exactness actually preserves}
\begin{proof}
An additive functor sends \(1_0=0_{0,0}\) to \(1_{F0}=0\), so
\(F0\) is a zero object. It preserves all biproduct equations,
including \(u_1p_1+u_2p_2=1\), hence takes a biproduct system to
one by its proved recognition criterion. Conversely, a functor
preserving zero and biproduct systems takes the canonical diagonal and
codiagonal to the corresponding canonical arrows, by their uniqueness.
Applying it to the diagonal/codiagonal expression for addition gives
\(F(f+g)=Ff+Fg\); it is additive.

Suppose a zero-preserving functor is left exact. On a biproduct, its
four projection/injection equations are preserved. In the target put
\(w=1-Fu_1Fp_1-Fu_2Fp_2\). Then \(Fp_1w=Fp_2w=0\).
The preserved left exact sequence with kernel \(Fu_2\) of \(Fp_1\)
gives \(w=Fu_2z\); composing with \(Fp_2\) gives \(z=0\).
Thus \(w=0\), so the image maps are a biproduct system and \(F\)
is additive. The opposite-category argument gives the right exact case.

For an additive functor, preservation of kernels is equivalent to
preservation of left exact sequences, because \(0\to K\to A\to B\)
is left exact precisely when its first map is a kernel of the second.
The dual statement gives cokernels and right exactness. If \(F\)
preserves short exact sequences, apply it to the kernel--image and
image--cokernel short sequences of an arbitrary \(f\). It preserves
the epic first factor and monic second factor. Consequently the first
preserved kernel is a kernel of \(Ff\), and the second preserved
cokernel is a cokernel of \(Ff\). Thus it preserves all kernels and
cokernels. Conversely, preservation of both carries canonical image
factorizations to image factorizations, since image is kernel of
cokernel. Therefore every equality \(\operatorname{Im}f=\operatorname{Ker}g\)
is preserved, and all finite exact sequences are preserved. Exactness
obviously includes the short sequences. This proves all equivalences.

For \((T,-)\), each induced map is postcomposition, a group
homomorphism by bilinearity. Given a kernel \(k:K\to A\) of \(g\),
an arrow \(x:T\to A\) is killed by \(g\) iff it has a unique
factorization through \(k\). This identifies \((T,K)\) with the
group kernel of \((T,g)\), so the functor is left exact. The concrete
integer countersequence in the next item proves that right exactness
does not follow.
\end{proof}
\end{annexitem}

\begin{annexitem}{annex:p10.ab.15}{thm:p10-ab-module-cases}{The concrete calculations and the nonzero connecting map}
\begin{proof}
In left \(R\)-modules the zero module is zero, and finite direct sums
are products and coproducts with their coordinate maps. For \(f:A\to B\),
the set \(\{a:f(a)=0\}\) is an \(R\)-submodule; its inclusion
has the kernel universal property. The quotient \(B/f[A]\) is an
\(R\)-module and its quotient map has the cokernel property: an
\(R\)-linear map killing \(f[A]\) is constant on its cosets and
factors uniquely. This is a calculation with actual module elements.
The map \(A/\ker f\to f[A]\), \([a]\mapsto f(a)\), is well
defined and linear, injective since equal values have difference in
\(\ker f\), and onto by definition of \(f[A]\). Its inverse is
linear by uniqueness of the representatives' classes.

A module mono is injective: a nonzero \(a\in\ker f\) would give
distinct maps \(R\to A\), \(r\mapsto ra\) and zero, equal after
\(f\). The converse is cancellation of functions. A module epi is
onto: otherwise the nonzero quotient map \(B\to B/f[A]\) and the
zero map would agree after \(f\), contradicting epicity. The converse
again follows from surjective-function cancellation. Thus every mono
is the kernel of its image quotient, and every epi is a cokernel of its
kernel, with the displayed module isomorphism supplying the universal
factorizations. All structural abelian axioms are established, including
the possible trivial category for the zero ring.

For the vector map \((x,y,z)\mapsto(x+y,z)\), its kernel consists
exactly of \((t,-t,0)\). The displayed section is linear and its
composite with the map is \(1_{k^2}\). The retraction
\((x,y,z)\mapsto-y\) satisfies \(ri=1\), \(rs=0\), and
\(ir+sp=1\), verifying split exactness and the actual inverse
\(k^3\cong k\oplus k^2\). The formulas work also in characteristic two.

For the integer diagram, commutativity holds because
\(n\circ n=n\circ n\) on its first square, and multiplication by
\(n\) induces zero modulo \(n\) on its second. In the construction
of \(\delta\), a class \([z]\in\operatorname{Ker}0=\mathbb Z/n\)
has a lift \(z\) in the upper middle \(\mathbb Z\); its image in
the lower middle is \(nz\), which is the image of \(z\) under the
lower left injection. Quotienting this \(z\) in
\(\operatorname{Cok}n\) gives \([z]\), independently of the lift.
Hence \(\delta\) is the identity, with the construction's positive
sign. The next induced map is zero, and the last is the identity:
the sequence reduces to
\[
 0\to0\to0\to\mathbb Z/n\xrightarrow{1}\mathbb Z/n
 \xrightarrow{0}\mathbb Z/n\xrightarrow{1}\mathbb Z/n\to0,
\]
which also verifies all its exactness positions directly in this example.
Every homomorphism \(\mathbb Z/n\to\mathbb Z\) is zero, since its
generator image is annihilated by \(n\) in a torsion-free group. There
can be no section of the integer quotient, so the short sequence is
nonsplit. Applying \((\mathbb Z/n,-)\) gives zero at both integer
terms but the nonzero endomorphism group \(\mathbb Z/n\) at the
quotient term. This explicitly fails right exactness.
\end{proof}
\end{annexitem}

\begin{annexitem}{annex:p10.ab.16}{thm:p10-ab-countercases}{Three failures outside the abelian axioms}
\begin{proof}
The inclusion of \(H=\{1,(12)\}\) in \(S_3\) is monic in groups.
Every kernel of a group homomorphism is normal, because
\(f(ghg^{-1})=f(g)f(h)f(g)^{-1}=1\) whenever \(f(h)=1\).
But conjugating \((12)\) by \((123)\) gives \((23)\notin H\).
Thus this mono is not a kernel, and A3 fails. No assertion about
epimorphisms of arbitrary groups is needed for this counterexample.

In pointed sets, the one-point object is zero. For the two two-point
objects \(\{*,a\},\{*,b\}\), their wedge has three points while
their product has four; the canonical map misses \((a,b)\), so cannot
be a biproduct. Let \(q:\{*,a,b\}\to\{*,c\}\) send both \(a,b\)
to \(c\). It is epic by surjective-function cancellation. If it were
a cokernel of \(g\), then \(qg=0\) and its basepoint fiber
\(q^{-1}(*)=\{*\}\) would force \(g=0\). The cokernel of a zero
arrow into this source is its identity, so cokernel uniqueness would
make \(q\) invertible, a contradiction. Hence A3* fails as well.

Finitely generated free abelian groups are closed under finite direct
sums, their hom sets are abelian groups, and composition is bilinear;
this is an additive category. Multiplication by two on \(\mathbb Z\)
is monic by injectivity. It is epic internally: for maps into any free
abelian group \(F\), \(r\circ2=s\circ2\) implies
\(2(r(1)-s(1))=0\), hence \(r=s\) by torsion-freeness.
It is not invertible, since an integer image of \(1\) cannot be a
half-integer. Its kernel is zero. Every map from \(\mathbb Z\) to
a free group killing two is zero, so its internal cokernel is zero,
even though its cokernel in all abelian groups is \(\mathbb Z/2\).
Therefore its internal coimage and image are both \(\mathbb Z\),
and the canonical comparison is the original multiplication by two,
not an isomorphism. This is a concrete failure of balance and normality
in an additive category, not an invocation of a desired comparison as
an axiom or a mistaken use of an object absent from the category.
\end{proof}
\end{annexitem}
